
\documentclass{article}
\usepackage{colortbl}
\usepackage{makecell}
\usepackage{multirow}
\usepackage{supertabular}

\begin{document}

\newcounter{utterance}

\centering \large Interaction Transcript for game `clean\_up', experiment `2\_hard\_3obj\_bg', episode 1 with mistral{-}large{-}3{-}azure.
\vspace{24pt}

{ \footnotesize  \setcounter{utterance}{1}
\setlength{\tabcolsep}{0pt}
\begin{supertabular}{c@{$\;$}|p{.15\linewidth}@{}p{.15\linewidth}p{.15\linewidth}p{.15\linewidth}p{.15\linewidth}p{.15\linewidth}}
    \# & $\;$A & \multicolumn{4}{c}{Game Master} & $\;\:$B\\
    \hline

    \theutterance \stepcounter{utterance}  
    & & \multicolumn{4}{p{0.6\linewidth}}{
        \cellcolor[rgb]{0.9,0.9,0.9}{
            \makecell[{{p{\linewidth}}}]{
                \texttt{\tiny{[P1$\langle$GM]}}
                \texttt{Аз съм вашият ръководител на играта и вие играете съвместна игра със следната мрежа като игрално поле:} \\
\\ 
\texttt{```} \\
\texttt{ 1234567} \\
\texttt{╔═══╤═╤═╗} \\
\texttt{║◌◌◌│◌│◌║ 1} \\
\texttt{║◌P◌└─┤◌║ 2} \\
\texttt{║◌◌◌C◌├─╢ 3} \\
\texttt{╟─┐◌┌─┤◌║ 4} \\
\texttt{║◌│◌│◌│◌║ 5} \\
\texttt{╟─┼─┼─┼─╢ 6} \\
\texttt{║◌│◌│◌│L║ 7} \\
\texttt{╚═╧═╧═╧═╝} \\
\\ 
\texttt{```} \\
\\ 
\texttt{* Горният ръб показва x{-}координати, които нарастват надясно, а десният ръб показва y{-}координати, които нарастват надолу.} \\
\texttt{* Следните обекти са разположени на случаен принцип върху вашата мрежа: 'C', 'L', 'P'.} \\
\\ 
\\ 
\texttt{Другият играч вижда вариант на игралното поле, където обектите са разположени на различни произволни места. Вие не можете да видите полето на другия играч, а той не може да види вашето.} \\
\\ 
\\ 
\texttt{**Цел:**} \\
\\ 
\texttt{И двамата играчи трябва да преместят обектите на своя фон така, че еднакви обекти да се окажат на едни и същи координати. Трябва да общувате с другия играч, за да се споразумеете за обща целева конфигурация.} \\
\\ 
\\ 
\texttt{**Правила:**} \\
\\ 
\texttt{* Във всеки ход можете да изпратите точно една от следните две команди:} \\
\texttt{ 1. `КАЖИ: <СЪОБЩЕНИЕ>`: за изпращане на съобщение (всичко до следващото прекъсване на реда) на другия играч. Аз ще го препратя на вашия партньор.} \\
\texttt{2. `ПРЕМЕСТИ: <ОБЕКТ>, (<X>, <Y>)`: за преместване на обект на нова позиция, където `<X>` е колоната, а `<Y>` е редът. Ще ви информирам дали преместването ви е успешно или не.} \\
\texttt{* Ако не се придържате към формата или изпратите няколко команди едновременно, ще трябва да ви накажа.} \\
\texttt{* Ако и двамата играчи натрупат повече от 8 наказания, и двамата губите играта.} \\
\texttt{* Изключително важно е да комуникирате с другия играч относно вашето целево състояние! Единственият начин, по който можете да предадете стратегията си на другия играч, е като използвате командата `КАЖИ: <СЪОБЩЕНИЕ>`!} \\
\\ 
\\ 
\texttt{**Преместване на обекти**} \\
\\ 
\texttt{* Можете да премествате обекти само в клетки в границите на мрежата. Целевата клетка трябва да е празна, т.е. да съдържа само символа '◌'.} \\
\texttt{* Ако се опитате да преместите обект на място, което не е празно, или се опитате да го преместите извън мрежата, ще трябва да ви накажа. Получавате още един опит.} \\
\texttt{* Преди да направите ход, проверете два пъти дали целевото място е празно и не съдържа никаква буква, рамка или линия!} \\
\\ 
\\ 
\texttt{**Край на играта**} \\
\\ 
\texttt{Ако смятате, че сте постигнали целта да подравните всички обекти, можете да помолите другия играч да завърши играта, като изпратите `КАЖИ: приключихте ли?`. Ако другият играч ви помоли да завършите играта и вие отговорите с `КАЖИ: приключихме!`, играта ще приключи.} \\
\\ 
\texttt{И двамата играчи печелят, ако играта приключи в рамките на 12 рунда, където един рунд се определя като двама играчи, всеки от които изпраща валидна команда.} \\
\\ 
\texttt{**Оценяване:**} \\
\\ 
\texttt{Колкото по{-}близо са еднаквите обекти в двете игрални полета, толкова повече точки получавате. Наказанията намаляват точките ви. Можете ли да подобрите рекорда?} \\
\\ 
\\ 
\texttt{Моля, изпратете съобщение на другия играч, за да започнете играта!} \\
            }
        }
    }
    & \\ \\

    \theutterance \stepcounter{utterance}  
    & \multicolumn{4}{p{0.6\linewidth}}{
        \cellcolor[rgb]{0.9,0.9,0.9}{
            \makecell[{{p{\linewidth}}}]{
                \texttt{\tiny{[P1$\rangle$GM]}}
                \texttt{Да започнем играта! Ето моето първо съобщение за другия играч, за да координираме действията си:} \\
\\ 
\texttt{```} \\
\texttt{КАЖИ: Здравей! Нека да съгласуваме координатите за всеки обект. Моите обекти са разположени така:} \\
\texttt{{-} P е на (3, 2)} \\
\texttt{{-} C е на (4, 3)} \\
\texttt{{-} L е на (7, 7)} \\
\\ 
\texttt{Може ли да споделиш къде са твоите обекти? След това можем да изберем общи координати за всеки от тях.} \\
\texttt{```} \\
\\ 
\texttt{След като получа информация за разположението на обектите на другия играч, ще предложа целеви координати, които да са удобни и за двама ни. Да се стремим към минимален брой ходове!} \\
            }
        }
    }
    & & \\ \\

    \theutterance \stepcounter{utterance}  
    & & & \multicolumn{2}{p{0.3\linewidth}}{
        \cellcolor[rgb]{0.9,0.9,0.9}{
            \makecell[{{p{\linewidth}}}]{
                \texttt{\tiny{[GM$|$GM]}}
                \texttt{Да започнем играта! Ето моето първо съобщение за другия играч, за да координираме действията си:} \\
\\ 
\texttt{```} \\
\texttt{КАЖИ: Здравей! Нека да съгласуваме координатите за всеки обект. Моите обекти са разположени така:} \\
\texttt{{-} P е на (3, 2)} \\
\texttt{{-} C е на (4, 3)} \\
\texttt{{-} L е на (7, 7)} \\
\\ 
\texttt{Може ли да споделиш къде са твоите обекти? След това можем да изберем общи координати за всеки от тях.} \\
\texttt{```} \\
\\ 
\texttt{След като получа информация за разположението на обектите на другия играч, ще предложа целеви координати, които да са удобни и за двама ни. Да се стремим към минимален брой ходове!} \\
            }
        }
    }
    & & \\ \\

    \theutterance \stepcounter{utterance}  
    & & & \multicolumn{2}{p{0.3\linewidth}}{
        \cellcolor[rgb]{0.9,0.9,0.9}{
            \makecell[{{p{\linewidth}}}]{
                \texttt{\tiny{[GM$|$GM]}}
                \texttt{Invalid format: head and tail are not empty} \\
\texttt{head: 'Да започнем играта! Ето моето първо съобщение за другия играч, за да координираме действията си:} \\
\\ 
\\ 
\texttt{'} \\
\texttt{tail: '} \\
\texttt{{-} P е на (3, 2)} \\
\texttt{{-} C е на (4, 3)} \\
\texttt{{-} L е на (7, 7)} \\
\\ 
\texttt{Може ли да споделиш къде са твоите обекти? След това можем да изберем общи координати за всеки от тях.} \\
\\ 
\\ 
\texttt{След като получа информация за разположението на обектите на другия играч, ще предложа целеви координати, които да са удобни и за двама ни. Да се стремим към минимален брой ходове!'} \\
            }
        }
    }
    & & \\ \\

    \theutterance \stepcounter{utterance}  
    & & \multicolumn{4}{p{0.6\linewidth}}{
        \cellcolor[rgb]{0.9,0.9,0.9}{
            \makecell[{{p{\linewidth}}}]{
                \texttt{\tiny{[P1$\langle$GM]}}
                \texttt{Наказание: Вашето съобщение не трябва да съдържа нищо преди или след командата!} \\
\\ 
\texttt{Уверете се, че вашият отговор съдържа само `КАЖИ: <СЪОБЩЕНИЕ>` или `ПРЕМЕСТИ: <ОБЕКТ>, (<X>, <Y>)` и нищо друго!} \\
\\ 
\texttt{Колективно сте натрупали 1 от 8 наказания.} \\
\texttt{Моля, опитайте отново!} \\
            }
        }
    }
    & \\ \\

    \theutterance \stepcounter{utterance}  
    & \multicolumn{4}{p{0.6\linewidth}}{
        \cellcolor[rgb]{0.9,0.9,0.9}{
            \makecell[{{p{\linewidth}}}]{
                \texttt{\tiny{[P1$\rangle$GM]}}
                \texttt{```} \\
\texttt{КАЖИ: Здравей! Как са разположени обектите ти? Моите са: P(3,2), C(4,3), L(7,7). Да изберем общи координати.} \\
\texttt{```} \\
            }
        }
    }
    & & \\ \\

    \theutterance \stepcounter{utterance}  
    & & & \multicolumn{2}{p{0.3\linewidth}}{
        \cellcolor[rgb]{0.9,0.9,0.9}{
            \makecell[{{p{\linewidth}}}]{
                \texttt{\tiny{[GM$|$GM]}}
                \texttt{```} \\
\texttt{КАЖИ: Здравей! Как са разположени обектите ти? Моите са: P(3,2), C(4,3), L(7,7). Да изберем общи координати.} \\
\texttt{```} \\
            }
        }
    }
    & & \\ \\

    \theutterance \stepcounter{utterance}  
    & & \multicolumn{4}{p{0.6\linewidth}}{
        \cellcolor[rgb]{0.9,0.9,0.9}{
            \makecell[{{p{\linewidth}}}]{
                \texttt{\tiny{[GM$\rangle$P2]}}
                \texttt{Аз съм вашият ръководител на играта и вие играете съвместна игра със следната мрежа като игрално поле:} \\
\\ 
\texttt{```} \\
\texttt{ 1234567} \\
\texttt{╔═══╤═╤═╗} \\
\texttt{║◌◌◌│◌│◌║ 1} \\
\texttt{║◌◌L└─┤◌║ 2} \\
\texttt{║◌◌◌◌◌├─╢ 3} \\
\texttt{╟─┐◌┌─┤◌║ 4} \\
\texttt{║C│◌│◌│◌║ 5} \\
\texttt{╟─┼─┼─┼─╢ 6} \\
\texttt{║P│◌│◌│◌║ 7} \\
\texttt{╚═╧═╧═╧═╝} \\
\\ 
\texttt{```} \\
\\ 
\texttt{* Горният ръб показва x{-}координати, които нарастват надясно, а десният ръб показва y{-}координати, които нарастват надолу.} \\
\texttt{* Следните обекти са разположени на случаен принцип върху вашата мрежа: 'C', 'L', 'P'.} \\
\\ 
\\ 
\texttt{Другият играч вижда вариант на игралното поле, където обектите са разположени на различни произволни места. Вие не можете да видите полето на другия играч, а той не може да види вашето.} \\
\\ 
\\ 
\texttt{**Цел:**} \\
\\ 
\texttt{И двамата играчи трябва да преместят обектите на своя фон така, че еднакви обекти да се окажат на едни и същи координати. Трябва да общувате с другия играч, за да се споразумеете за обща целева конфигурация.} \\
\\ 
\\ 
\texttt{**Правила:**} \\
\\ 
\texttt{* Във всеки ход можете да изпратите точно една от следните две команди:} \\
\texttt{ 1. `КАЖИ: <СЪОБЩЕНИЕ>`: за изпращане на съобщение (всичко до следващото прекъсване на реда) на другия играч. Аз ще го препратя на вашия партньор.} \\
\texttt{2. `ПРЕМЕСТИ: <ОБЕКТ>, (<X>, <Y>)`: за преместване на обект на нова позиция, където `<X>` е колоната, а `<Y>` е редът. Ще ви информирам дали преместването ви е успешно или не.} \\
\texttt{* Ако не се придържате към формата или изпратите няколко команди едновременно, ще трябва да ви накажа.} \\
\texttt{* Ако и двамата играчи натрупат повече от 8 наказания, и двамата губите играта.} \\
\texttt{* Изключително важно е да комуникирате с другия играч относно вашето целево състояние! Единственият начин, по който можете да предадете стратегията си на другия играч, е като използвате командата `КАЖИ: <СЪОБЩЕНИЕ>`!} \\
\\ 
\\ 
\texttt{**Преместване на обекти**} \\
\\ 
\texttt{* Можете да премествате обекти само в клетки в границите на мрежата. Целевата клетка трябва да е празна, т.е. да съдържа само символа '◌'.} \\
\texttt{* Ако се опитате да преместите обект на място, което не е празно, или се опитате да го преместите извън мрежата, ще трябва да ви накажа. Получавате още един опит.} \\
\texttt{* Преди да направите ход, проверете два пъти дали целевото място е празно и не съдържа никаква буква, рамка или линия!} \\
\\ 
\\ 
\texttt{**Край на играта**} \\
\\ 
\texttt{Ако смятате, че сте постигнали целта да подравните всички обекти, можете да помолите другия играч да завърши играта, като изпратите `КАЖИ: приключихте ли?`. Ако другият играч ви помоли да завършите играта и вие отговорите с `КАЖИ: приключихме!`, играта ще приключи.} \\
\\ 
\texttt{И двамата играчи печелят, ако играта приключи в рамките на 12 рунда, където един рунд се определя като двама играчи, всеки от които изпраща валидна команда.} \\
\\ 
\texttt{**Оценяване:**} \\
\\ 
\texttt{Колкото по{-}близо са еднаквите обекти в двете игрални полета, толкова повече точки получавате. Наказанията намаляват точките ви. Можете ли да подобрите рекорда?} \\
\\ 
\texttt{Другият играч започна играта, като изпрати това съобщение:} \\
\\ 
\texttt{"Здравей! Как са разположени обектите ти? Моите са: P(3,2), C(4,3), L(7,7). Да изберем общи координати."} \\
\\ 
\texttt{Каква е първата ви команда?} \\
            }
        }
    }
    & \\ \\

    \theutterance \stepcounter{utterance}  
    & & & \multicolumn{4}{p{0.6\linewidth}}{
        \cellcolor[rgb]{0.9,0.9,0.9}{
            \makecell[{{p{\linewidth}}}]{
                \texttt{\tiny{[GM$\langle$P2]}}
                \texttt{Отговорът на другия играч показва неговите координати за обектите. За да започнем да се синхронизираме, първо трябва да споделим местоположението на нашите обекти и след това да предложим общи целеви координати. Ето каква може да бъде първата ви команда:} \\
\\ 
\texttt{```} \\
\texttt{КАЖИ: Моите обекти са разположени така: C(2,5), L(3,2), P(2,7). Предлагам да изберем следните общи координати: C(4,5), L(3,3), P(5,7). Какво мислиш?} \\
\texttt{```} \\
\\ 
\texttt{Това съобщение:} \\
\texttt{1. Споделя вашите текущи координати.} \\
\texttt{2. Предлага общи целеви координати, които са възможни за постигане (празни клетки на вашето поле).} \\
\texttt{3. Поканва другия играч да коментира или да предложи алтернативи.} \\
            }
        }
    }
     \\ \\

    \theutterance \stepcounter{utterance}  
    & & & \multicolumn{2}{p{0.3\linewidth}}{
        \cellcolor[rgb]{0.9,0.9,0.9}{
            \makecell[{{p{\linewidth}}}]{
                \texttt{\tiny{[GM$|$GM]}}
                \texttt{Отговорът на другия играч показва неговите координати за обектите. За да започнем да се синхронизираме, първо трябва да споделим местоположението на нашите обекти и след това да предложим общи целеви координати. Ето каква може да бъде първата ви команда:} \\
\\ 
\texttt{```} \\
\texttt{КАЖИ: Моите обекти са разположени така: C(2,5), L(3,2), P(2,7). Предлагам да изберем следните общи координати: C(4,5), L(3,3), P(5,7). Какво мислиш?} \\
\texttt{```} \\
\\ 
\texttt{Това съобщение:} \\
\texttt{1. Споделя вашите текущи координати.} \\
\texttt{2. Предлага общи целеви координати, които са възможни за постигане (празни клетки на вашето поле).} \\
\texttt{3. Поканва другия играч да коментира или да предложи алтернативи.} \\
            }
        }
    }
    & & \\ \\

    \theutterance \stepcounter{utterance}  
    & & & \multicolumn{2}{p{0.3\linewidth}}{
        \cellcolor[rgb]{0.9,0.9,0.9}{
            \makecell[{{p{\linewidth}}}]{
                \texttt{\tiny{[GM$|$GM]}}
                \texttt{Invalid format: head and tail are not empty} \\
\texttt{head: 'Отговорът на другия играч показва неговите координати за обектите. За да започнем да се синхронизираме, първо трябва да споделим местоположението на нашите обекти и след това да предложим общи целеви координати. Ето каква може да бъде първата ви команда:} \\
\\ 
\\ 
\texttt{'} \\
\texttt{tail: '} \\
\\ 
\\ 
\texttt{Това съобщение:} \\
\texttt{1. Споделя вашите текущи координати.} \\
\texttt{2. Предлага общи целеви координати, които са възможни за постигане (празни клетки на вашето поле).} \\
\texttt{3. Поканва другия играч да коментира или да предложи алтернативи.'} \\
            }
        }
    }
    & & \\ \\

    \theutterance \stepcounter{utterance}  
    & & \multicolumn{4}{p{0.6\linewidth}}{
        \cellcolor[rgb]{0.9,0.9,0.9}{
            \makecell[{{p{\linewidth}}}]{
                \texttt{\tiny{[GM$\rangle$P2]}}
                \texttt{Наказание: Вашето съобщение не трябва да съдържа нищо преди или след командата!} \\
\\ 
\texttt{Уверете се, че вашият отговор съдържа само `КАЖИ: <СЪОБЩЕНИЕ>` или `ПРЕМЕСТИ: <ОБЕКТ>, (<X>, <Y>)` и нищо друго!} \\
\\ 
\texttt{Колективно сте натрупали 2 от 8 наказания.} \\
\texttt{Моля, опитайте отново!} \\
            }
        }
    }
    & \\ \\

    \theutterance \stepcounter{utterance}  
    & & & \multicolumn{4}{p{0.6\linewidth}}{
        \cellcolor[rgb]{0.9,0.9,0.9}{
            \makecell[{{p{\linewidth}}}]{
                \texttt{\tiny{[GM$\langle$P2]}}
                \texttt{Разбрано! Ето коректната команда:} \\
\\ 
\texttt{```} \\
\texttt{КАЖИ: Моите обекти: C(2,5), L(3,2), P(2,7). Да изберем C(4,5), L(3,3), P(5,7)?} \\
\texttt{```} \\
            }
        }
    }
     \\ \\

    \theutterance \stepcounter{utterance}  
    & & & \multicolumn{2}{p{0.3\linewidth}}{
        \cellcolor[rgb]{0.9,0.9,0.9}{
            \makecell[{{p{\linewidth}}}]{
                \texttt{\tiny{[GM$|$GM]}}
                \texttt{Разбрано! Ето коректната команда:} \\
\\ 
\texttt{```} \\
\texttt{КАЖИ: Моите обекти: C(2,5), L(3,2), P(2,7). Да изберем C(4,5), L(3,3), P(5,7)?} \\
\texttt{```} \\
            }
        }
    }
    & & \\ \\

    \theutterance \stepcounter{utterance}  
    & & & \multicolumn{2}{p{0.3\linewidth}}{
        \cellcolor[rgb]{0.9,0.9,0.9}{
            \makecell[{{p{\linewidth}}}]{
                \texttt{\tiny{[GM$|$GM]}}
                \texttt{Invalid format: head is not empty: 'Разбрано! Ето коректната команда:} \\
\\ 
\\ 
\texttt{'} \\
            }
        }
    }
    & & \\ \\

    \theutterance \stepcounter{utterance}  
    & & \multicolumn{4}{p{0.6\linewidth}}{
        \cellcolor[rgb]{0.9,0.9,0.9}{
            \makecell[{{p{\linewidth}}}]{
                \texttt{\tiny{[GM$\rangle$P2]}}
                \texttt{Наказание: Вашето съобщение не трябва да съдържа нищо преди командата!} \\
\\ 
\texttt{Уверете се, че вашият отговор съдържа само `КАЖИ: <СЪОБЩЕНИЕ>` или `ПРЕМЕСТИ: <ОБЕКТ>, (<X>, <Y>)` и нищо друго!} \\
\\ 
\texttt{Колективно сте натрупали 3 от 8 наказания.} \\
\texttt{Моля, опитайте отново!} \\
            }
        }
    }
    & \\ \\

    \theutterance \stepcounter{utterance}  
    & & & \multicolumn{4}{p{0.6\linewidth}}{
        \cellcolor[rgb]{0.9,0.9,0.9}{
            \makecell[{{p{\linewidth}}}]{
                \texttt{\tiny{[GM$\langle$P2]}}
                \texttt{Ясно! Ето правилната команда без никакви допълнения:} \\
\\ 
\texttt{```} \\
\texttt{КАЖИ: C(2,5), L(3,2), P(2,7). Предлагам C(4,5), L(3,3), P(5,7)} \\
\texttt{```} \\
            }
        }
    }
     \\ \\

    \theutterance \stepcounter{utterance}  
    & & & \multicolumn{2}{p{0.3\linewidth}}{
        \cellcolor[rgb]{0.9,0.9,0.9}{
            \makecell[{{p{\linewidth}}}]{
                \texttt{\tiny{[GM$|$GM]}}
                \texttt{Ясно! Ето правилната команда без никакви допълнения:} \\
\\ 
\texttt{```} \\
\texttt{КАЖИ: C(2,5), L(3,2), P(2,7). Предлагам C(4,5), L(3,3), P(5,7)} \\
\texttt{```} \\
            }
        }
    }
    & & \\ \\

    \theutterance \stepcounter{utterance}  
    & & & \multicolumn{2}{p{0.3\linewidth}}{
        \cellcolor[rgb]{0.9,0.9,0.9}{
            \makecell[{{p{\linewidth}}}]{
                \texttt{\tiny{[GM$|$GM]}}
                \texttt{Invalid format: head is not empty: 'Ясно! Ето правилната команда без никакви допълнения:} \\
\\ 
\\ 
\texttt{'} \\
            }
        }
    }
    & & \\ \\

    \theutterance \stepcounter{utterance}  
    & & \multicolumn{4}{p{0.6\linewidth}}{
        \cellcolor[rgb]{0.9,0.9,0.9}{
            \makecell[{{p{\linewidth}}}]{
                \texttt{\tiny{[GM$\rangle$P2]}}
                \texttt{Наказание: Вашето съобщение не трябва да съдържа нищо преди командата!} \\
\\ 
\texttt{Уверете се, че вашият отговор съдържа само `КАЖИ: <СЪОБЩЕНИЕ>` или `ПРЕМЕСТИ: <ОБЕКТ>, (<X>, <Y>)` и нищо друго!} \\
\\ 
\texttt{Колективно сте натрупали 4 от 8 наказания.} \\
\texttt{Моля, опитайте отново!} \\
            }
        }
    }
    & \\ \\

    \theutterance \stepcounter{utterance}  
    & & & \multicolumn{4}{p{0.6\linewidth}}{
        \cellcolor[rgb]{0.9,0.9,0.9}{
            \makecell[{{p{\linewidth}}}]{
                \texttt{\tiny{[GM$\langle$P2]}}
                \texttt{Съжалявам за грешките. Ето абсолютно чистата команда:} \\
\\ 
\texttt{```} \\
\texttt{КАЖИ: Моите обекти са C(2,5), L(3,2), P(2,7). Какви са вашите предложения за общи координати?} \\
\texttt{```} \\
            }
        }
    }
     \\ \\

    \theutterance \stepcounter{utterance}  
    & & & \multicolumn{2}{p{0.3\linewidth}}{
        \cellcolor[rgb]{0.9,0.9,0.9}{
            \makecell[{{p{\linewidth}}}]{
                \texttt{\tiny{[GM$|$GM]}}
                \texttt{Съжалявам за грешките. Ето абсолютно чистата команда:} \\
\\ 
\texttt{```} \\
\texttt{КАЖИ: Моите обекти са C(2,5), L(3,2), P(2,7). Какви са вашите предложения за общи координати?} \\
\texttt{```} \\
            }
        }
    }
    & & \\ \\

    \theutterance \stepcounter{utterance}  
    & & & \multicolumn{2}{p{0.3\linewidth}}{
        \cellcolor[rgb]{0.9,0.9,0.9}{
            \makecell[{{p{\linewidth}}}]{
                \texttt{\tiny{[GM$|$GM]}}
                \texttt{Invalid format: head is not empty: 'Съжалявам за грешките. Ето абсолютно чистата команда:} \\
\\ 
\\ 
\texttt{'} \\
            }
        }
    }
    & & \\ \\

    \theutterance \stepcounter{utterance}  
    & & \multicolumn{4}{p{0.6\linewidth}}{
        \cellcolor[rgb]{0.9,0.9,0.9}{
            \makecell[{{p{\linewidth}}}]{
                \texttt{\tiny{[GM$\rangle$P2]}}
                \texttt{Наказание: Вашето съобщение не трябва да съдържа нищо преди командата!} \\
\\ 
\texttt{Уверете се, че вашият отговор съдържа само `КАЖИ: <СЪОБЩЕНИЕ>` или `ПРЕМЕСТИ: <ОБЕКТ>, (<X>, <Y>)` и нищо друго!} \\
\\ 
\texttt{Колективно сте натрупали 5 от 8 наказания.} \\
\texttt{Моля, опитайте отново!} \\
            }
        }
    }
    & \\ \\

    \theutterance \stepcounter{utterance}  
    & & & \multicolumn{4}{p{0.6\linewidth}}{
        \cellcolor[rgb]{0.9,0.9,0.9}{
            \makecell[{{p{\linewidth}}}]{
                \texttt{\tiny{[GM$\langle$P2]}}
                \texttt{Разбрах! Ето точната команда:} \\
\\ 
\texttt{```} \\
\texttt{КАЖИ: C(2,5) L(3,2) P(2,7)} \\
\texttt{```} \\
            }
        }
    }
     \\ \\

    \theutterance \stepcounter{utterance}  
    & & & \multicolumn{2}{p{0.3\linewidth}}{
        \cellcolor[rgb]{0.9,0.9,0.9}{
            \makecell[{{p{\linewidth}}}]{
                \texttt{\tiny{[GM$|$GM]}}
                \texttt{Разбрах! Ето точната команда:} \\
\\ 
\texttt{```} \\
\texttt{КАЖИ: C(2,5) L(3,2) P(2,7)} \\
\texttt{```} \\
            }
        }
    }
    & & \\ \\

    \theutterance \stepcounter{utterance}  
    & & & \multicolumn{2}{p{0.3\linewidth}}{
        \cellcolor[rgb]{0.9,0.9,0.9}{
            \makecell[{{p{\linewidth}}}]{
                \texttt{\tiny{[GM$|$GM]}}
                \texttt{Invalid format: head is not empty: 'Разбрах! Ето точната команда:} \\
\\ 
\\ 
\texttt{'} \\
            }
        }
    }
    & & \\ \\

    \theutterance \stepcounter{utterance}  
    & & \multicolumn{4}{p{0.6\linewidth}}{
        \cellcolor[rgb]{0.9,0.9,0.9}{
            \makecell[{{p{\linewidth}}}]{
                \texttt{\tiny{[GM$\rangle$P2]}}
                \texttt{Наказание: Вашето съобщение не трябва да съдържа нищо преди командата!} \\
\\ 
\texttt{Уверете се, че вашият отговор съдържа само `КАЖИ: <СЪОБЩЕНИЕ>` или `ПРЕМЕСТИ: <ОБЕКТ>, (<X>, <Y>)` и нищо друго!} \\
\\ 
\texttt{Колективно сте натрупали 6 от 8 наказания.} \\
\texttt{Моля, опитайте отново!} \\
            }
        }
    }
    & \\ \\

    \theutterance \stepcounter{utterance}  
    & & & \multicolumn{4}{p{0.6\linewidth}}{
        \cellcolor[rgb]{0.9,0.9,0.9}{
            \makecell[{{p{\linewidth}}}]{
                \texttt{\tiny{[GM$\langle$P2]}}
                \texttt{Ясно! Ето съвсем минималистичната версия:} \\
\\ 
\texttt{```} \\
\texttt{КАЖИ: C2,5 L3,2 P2,7} \\
\texttt{```} \\
            }
        }
    }
     \\ \\

    \theutterance \stepcounter{utterance}  
    & & & \multicolumn{2}{p{0.3\linewidth}}{
        \cellcolor[rgb]{0.9,0.9,0.9}{
            \makecell[{{p{\linewidth}}}]{
                \texttt{\tiny{[GM$|$GM]}}
                \texttt{Ясно! Ето съвсем минималистичната версия:} \\
\\ 
\texttt{```} \\
\texttt{КАЖИ: C2,5 L3,2 P2,7} \\
\texttt{```} \\
            }
        }
    }
    & & \\ \\

    \theutterance \stepcounter{utterance}  
    & & & \multicolumn{2}{p{0.3\linewidth}}{
        \cellcolor[rgb]{0.9,0.9,0.9}{
            \makecell[{{p{\linewidth}}}]{
                \texttt{\tiny{[GM$|$GM]}}
                \texttt{Invalid format: head is not empty: 'Ясно! Ето съвсем минималистичната версия:} \\
\\ 
\\ 
\texttt{'} \\
            }
        }
    }
    & & \\ \\

    \theutterance \stepcounter{utterance}  
    & & \multicolumn{4}{p{0.6\linewidth}}{
        \cellcolor[rgb]{0.9,0.9,0.9}{
            \makecell[{{p{\linewidth}}}]{
                \texttt{\tiny{[GM$\rangle$P2]}}
                \texttt{Наказание: Вашето съобщение не трябва да съдържа нищо преди командата!} \\
\\ 
\texttt{Уверете се, че вашият отговор съдържа само `КАЖИ: <СЪОБЩЕНИЕ>` или `ПРЕМЕСТИ: <ОБЕКТ>, (<X>, <Y>)` и нищо друго!} \\
\\ 
\texttt{Колективно сте натрупали 7 от 8 наказания.} \\
\texttt{Моля, опитайте отново!} \\
            }
        }
    }
    & \\ \\

    \theutterance \stepcounter{utterance}  
    & & & \multicolumn{4}{p{0.6\linewidth}}{
        \cellcolor[rgb]{0.9,0.9,0.9}{
            \makecell[{{p{\linewidth}}}]{
                \texttt{\tiny{[GM$\langle$P2]}}
                \texttt{Съжалявам за грешките. Ето абсолютно коректната команда:} \\
\\ 
\texttt{```} \\
\texttt{КАЖИ: Моите координати: C(2,5) L(3,2) P(2,7)} \\
\texttt{```} \\
\\ 
\texttt{Ако и това не е прието, моля уточнете какъв точно формат се изисква за съобщението, за да избегнем още наказания. Благодаря!} \\
            }
        }
    }
     \\ \\

    \theutterance \stepcounter{utterance}  
    & & & \multicolumn{2}{p{0.3\linewidth}}{
        \cellcolor[rgb]{0.9,0.9,0.9}{
            \makecell[{{p{\linewidth}}}]{
                \texttt{\tiny{[GM$|$GM]}}
                \texttt{Съжалявам за грешките. Ето абсолютно коректната команда:} \\
\\ 
\texttt{```} \\
\texttt{КАЖИ: Моите координати: C(2,5) L(3,2) P(2,7)} \\
\texttt{```} \\
\\ 
\texttt{Ако и това не е прието, моля уточнете какъв точно формат се изисква за съобщението, за да избегнем още наказания. Благодаря!} \\
            }
        }
    }
    & & \\ \\

    \theutterance \stepcounter{utterance}  
    & & & \multicolumn{2}{p{0.3\linewidth}}{
        \cellcolor[rgb]{0.9,0.9,0.9}{
            \makecell[{{p{\linewidth}}}]{
                \texttt{\tiny{[GM$|$GM]}}
                \texttt{Invalid format: head and tail are not empty} \\
\texttt{head: 'Съжалявам за грешките. Ето абсолютно коректната команда:} \\
\\ 
\\ 
\texttt{'} \\
\texttt{tail: '} \\
\\ 
\\ 
\texttt{Ако и това не е прието, моля уточнете какъв точно формат се изисква за съобщението, за да избегнем още наказания. Благодаря!'} \\
            }
        }
    }
    & & \\ \\

    \theutterance \stepcounter{utterance}  
    & & \multicolumn{4}{p{0.6\linewidth}}{
        \cellcolor[rgb]{0.9,0.9,0.9}{
            \makecell[{{p{\linewidth}}}]{
                \texttt{\tiny{[GM$\rangle$P2]}}
                \texttt{Наказание: Вашето съобщение не трябва да съдържа нищо преди или след командата!} \\
\\ 
\texttt{Уверете се, че вашият отговор съдържа само `КАЖИ: <СЪОБЩЕНИЕ>` или `ПРЕМЕСТИ: <ОБЕКТ>, (<X>, <Y>)` и нищо друго!} \\
\\ 
\texttt{Колективно сте натрупали 8 от 8 наказания.} \\
\texttt{Моля, опитайте отново!} \\
            }
        }
    }
    & \\ \\

    \theutterance \stepcounter{utterance}  
    & & & \multicolumn{4}{p{0.6\linewidth}}{
        \cellcolor[rgb]{0.9,0.9,0.9}{
            \makecell[{{p{\linewidth}}}]{
                \texttt{\tiny{[GM$\langle$P2]}}
                \texttt{Разбирам изискванията сега. Ето перфектно форматираната команда:} \\
\\ 
\texttt{```} \\
\texttt{КАЖИ:C(2,5)L(3,2)P(2,7)} \\
\texttt{```} \\
\\ 
\texttt{Това е:} \\
\texttt{1. Само командата `КАЖИ:`} \\
\texttt{2. Без интервали след двоеточието} \\
\texttt{3. Без никакви допълнителни думи или символи} \\
\texttt{4. Само координатите на обектите в последователен формат} \\
\\ 
\texttt{Ако и това не е прието, моля да ми посочите точния формат, който се изисква. Благодаря за търпението!} \\
            }
        }
    }
     \\ \\

    \theutterance \stepcounter{utterance}  
    & & & \multicolumn{2}{p{0.3\linewidth}}{
        \cellcolor[rgb]{0.9,0.9,0.9}{
            \makecell[{{p{\linewidth}}}]{
                \texttt{\tiny{[GM$|$GM]}}
                \texttt{Разбирам изискванията сега. Ето перфектно форматираната команда:} \\
\\ 
\texttt{```} \\
\texttt{КАЖИ:C(2,5)L(3,2)P(2,7)} \\
\texttt{```} \\
\\ 
\texttt{Това е:} \\
\texttt{1. Само командата `КАЖИ:`} \\
\texttt{2. Без интервали след двоеточието} \\
\texttt{3. Без никакви допълнителни думи или символи} \\
\texttt{4. Само координатите на обектите в последователен формат} \\
\\ 
\texttt{Ако и това не е прието, моля да ми посочите точния формат, който се изисква. Благодаря за търпението!} \\
            }
        }
    }
    & & \\ \\

    \theutterance \stepcounter{utterance}  
    & & & \multicolumn{2}{p{0.3\linewidth}}{
        \cellcolor[rgb]{0.9,0.9,0.9}{
            \makecell[{{p{\linewidth}}}]{
                \texttt{\tiny{[GM$|$GM]}}
                \texttt{Invalid response format} \\
            }
        }
    }
    & & \\ \\

    \theutterance \stepcounter{utterance}  
    & & & \multicolumn{2}{p{0.3\linewidth}}{
        \cellcolor[rgb]{0.9,0.9,0.9}{
            \makecell[{{p{\linewidth}}}]{
                \texttt{\tiny{[GM$|$GM]}}
                \texttt{Maximum number of penalties exceeded} \\
            }
        }
    }
    & & \\ \\

    \theutterance \stepcounter{utterance}  
    & & & \multicolumn{2}{p{0.3\linewidth}}{
        \cellcolor[rgb]{0.9,0.9,0.9}{
            \makecell[{{p{\linewidth}}}]{
                \texttt{\tiny{[GM$|$GM]}}
                \texttt{Initial states:} \\
\texttt{```} \\
\texttt{Player 1:} \\
\texttt{ 1234567} \\
\texttt{╔═══╤═╤═╗} \\
\texttt{║◌◌◌│◌│◌║ 1} \\
\texttt{║◌P◌└─┤◌║ 2} \\
\texttt{║◌◌◌C◌├─╢ 3} \\
\texttt{╟─┐◌┌─┤◌║ 4} \\
\texttt{║◌│◌│◌│◌║ 5} \\
\texttt{╟─┼─┼─┼─╢ 6} \\
\texttt{║◌│◌│◌│L║ 7} \\
\texttt{╚═╧═╧═╧═╝} \\
\\ 
\\ 
\texttt{Player 2:} \\
\texttt{ 1234567} \\
\texttt{╔═══╤═╤═╗} \\
\texttt{║◌◌◌│◌│◌║ 1} \\
\texttt{║◌◌L└─┤◌║ 2} \\
\texttt{║◌◌◌◌◌├─╢ 3} \\
\texttt{╟─┐◌┌─┤◌║ 4} \\
\texttt{║C│◌│◌│◌║ 5} \\
\texttt{╟─┼─┼─┼─╢ 6} \\
\texttt{║P│◌│◌│◌║ 7} \\
\texttt{╚═╧═╧═╧═╝} \\
\\ 
\texttt{```} \\
            }
        }
    }
    & & \\ \\

    \theutterance \stepcounter{utterance}  
    & & & \multicolumn{2}{p{0.3\linewidth}}{
        \cellcolor[rgb]{0.9,0.9,0.9}{
            \makecell[{{p{\linewidth}}}]{
                \texttt{\tiny{[GM$|$GM]}}
                \texttt{End states:} \\
\texttt{```} \\
\texttt{Player 1:} \\
\texttt{ 1234567} \\
\texttt{╔═══╤═╤═╗} \\
\texttt{║◌◌◌│◌│◌║ 1} \\
\texttt{║◌P◌└─┤◌║ 2} \\
\texttt{║◌◌◌C◌├─╢ 3} \\
\texttt{╟─┐◌┌─┤◌║ 4} \\
\texttt{║◌│◌│◌│◌║ 5} \\
\texttt{╟─┼─┼─┼─╢ 6} \\
\texttt{║◌│◌│◌│L║ 7} \\
\texttt{╚═╧═╧═╧═╝} \\
\\ 
\\ 
\texttt{Player 2:} \\
\texttt{ 1234567} \\
\texttt{╔═══╤═╤═╗} \\
\texttt{║◌◌◌│◌│◌║ 1} \\
\texttt{║◌◌L└─┤◌║ 2} \\
\texttt{║◌◌◌◌◌├─╢ 3} \\
\texttt{╟─┐◌┌─┤◌║ 4} \\
\texttt{║C│◌│◌│◌║ 5} \\
\texttt{╟─┼─┼─┼─╢ 6} \\
\texttt{║P│◌│◌│◌║ 7} \\
\texttt{╚═╧═╧═╧═╝} \\
\\ 
\texttt{```} \\
            }
        }
    }
    & & \\ \\

    \theutterance \stepcounter{utterance}  
    & & & \multicolumn{2}{p{0.3\linewidth}}{
        \cellcolor[rgb]{0.9,0.9,0.9}{
            \makecell[{{p{\linewidth}}}]{
                \texttt{\tiny{[GM$|$GM]}}
                \texttt{* success: False} \\
\texttt{* lose: False} \\
\texttt{* aborted: True} \\
\texttt{{-}{-}{-}{-}{-}{-}{-}} \\
\texttt{* Shifts: 0.00} \\
\texttt{* Max Shifts: 4.00} \\
\texttt{* Min Shifts: 2.00} \\
\texttt{* End Distance Sum: 15.11} \\
\texttt{* Init Distance Sum: 15.11} \\
\texttt{* Expected Distance Sum: 12.57} \\
\texttt{* Penalties: 9.00} \\
\texttt{* Max Penalties: 8.00} \\
\texttt{* Rounds: 0.00} \\
\texttt{* Max Rounds: 12.00} \\
\texttt{* Object Count: 3.00} \\
            }
        }
    }
    & & \\ \\

\end{supertabular}
}

\end{document}
